% Setup Table S1 -- dataset comparison. \input{} from report-draft.tex Ch6.
\begin{table}[tbp]
\centering
\small
\begin{tabular}{l l l}
\toprule
 & \textbf{PDB} & \textbf{AFDB (Swiss-Prot)} \\
\midrule
Origin                  & Experimental (X-ray, cryo-EM) & Predicted (AlphaFold2, v6) \\
Per-residue confidence  & B-factor & pLDDT \\
Source pool             & 194{,}090 structures & 550{,}122 structures \\
Train / val / test chains & 425{,}100 / 5{,}000 / 500 & 229{,}670 / 4{,}521 / 235 \\
Length filter           & 50--512 res ($n{=}256,128$ subsets) & $\le 256$ res \\
CATH coverage (train)   & 44.2\% & 61.8\% \\
Dominant CATH class     & $\alpha/\beta$ (55\%) & $\alpha/\beta$ (57\%) \\
Split                   & random, 98/1.9/0.1 & random, 98/1.9/0.1 \\
Val$\leftrightarrow$train, $\ge$30\% seq.\ id. & 98.3\% & 92.2\% \\
Val$\leftrightarrow$train, identical sequence  & 79.0\% & 29.6\% \\
\bottomrule
\end{tabular}
\caption{\textbf{Our two regimes bracket the data problem --- experimental-but-scarce PDB versus synthetic-but-abundant AFDB --- and both have heavily leaked splits.} AFDB, being machine-generated, is also more in-silico-designable. Neither split is sequence-clustered (79\% of PDB validation chains share a byte-identical sequence with training), so throughout we read rank-order and $\Delta$-from-baseline rather than absolute in-house scores (\S\ref{ch:evaluation}).}
\label{tab:setup-datasets}
\end{table}
